\input{sol_pre.tex}
\input{q_2024-12-01_1.tex}
\begin{document}
\input{sol_head.tex}

\Savol{1}
$36=4\cdot9$. Raqamlar yiğindisi: $7+a+b$ $\vdots$ $9$ $\Rightarrow$ $a+b=2$ yoki $11$; \ $\overline{ab}$ $\vdots$ $4$.\\
$a+b=2$: $20$; \ $a+b=11$: $56$, $92$. \ Jami $3$ ta.
\Javob{A}

\Savol{2}
$4{,}(36)=4\frac4{11}=\frac{48}{11}$, \ $3{,}1(9)=3{,}2=\frac{16}5$, \ $\dfrac1{3\frac17}=\dfrac7{22}$:
\[\frac{48}{11}\cdot\frac5{16}+\frac7{22}=\frac{30+7}{22}=1\frac{15}{22}.\]
\Javob{B}

\Savol{3}
$v+u=\dfrac S{80}$, \ $v-u=\dfrac{S/2}{60}=\dfrac S{120}$ $\Rightarrow$ $\dfrac{v+u}{v-u}=\dfrac32$ $\Rightarrow$ $v=5u$.
\Javob{D}

\Savol{4}
$15\,000\,000\cdot0{,}4\cdot0{,}3=1\,800\,000$ söm.
\Javob{A}

\Savol{5}
$s=10^z$: \ $x^2+y^2=10s$, \ $xy=\dfrac{s^2-10s}2$.
\[x^3+y^3=s^3-3xy\,s=s^3-\frac32\left(s^3-10s^2\right)=-\frac12s^3+15s^2.\]
$a=-\frac12$, $b=15$: \ $a+b=14{,}5$.
\Javob{D}

\Savol{6}
$\dfrac1{5+2\sqrt6}=5-2\sqrt6$: \ $5-2\sqrt6-5+2\sqrt6=0$.
\Javob{C}

\Savol{7}
$21-12\sqrt3=\left(2\sqrt3-3\right)^2$, \ $\left|\sqrt3-4\right|=4-\sqrt3$:
\[\left(2\sqrt3-3\right)\left(4-\sqrt3\right)=8\sqrt3-6-12+3\sqrt3=11\sqrt3-18.\]
\Javob{A}

\Savol{8}
$a_1,a_3,\dots,a_{2n-1}$ — $n$ ta had, örta qiymati $a_n$: \ $n\,a_n=96$. \ $a_{n-2}+a_{n+2}=2a_n=12$ $\Rightarrow$ $a_n=6$, $n=16$.
\Javob{D}

\Savol{9}
$\dfrac{b_1}{1-q}=22{,}5$, \ $b_1(1+q)=20$ $\Rightarrow$ $22{,}5\left(1-q^2\right)=20$, \ $q^2=\frac19$, $q=\frac13$, $b_1=15$.\\
$b_4=15\cdot\frac1{27}=\frac59$.
\Javob{D}

\Savol{10}
$a-b=1$: \ $a^3-b^3=a^2+ab+b^2$.
\[a^2+ab+b^2+a^2+b^2+2-5ab=2(a-b)^2+2=4.\]
\Javob{A}

\Savol{11}
Surat: $xy^2+x-x^2y-y=(x-y)(1-xy)$. \ Maxraj: $xy(x-y)+(x-y)(x+y)=(x-y)(xy+x+y)$.
\[\frac{1-xy}{xy+x+y}=\frac{1+1}{-1+3}=1.\]
\Javob{C}

\Savol{12}
$12345''=205'\,45''=3^\circ25'45''$.
\Javob{A}

\Savol{13}
Surat: $2\sin^2x\cos^2x\left(\cos^2x-\sin^2x\right)=\frac12\sin^22x\cos2x$.\\
Maxraj: $1-\cos^24x\left(\sin^24x+\cos^24x\right)=\sin^24x=4\sin^22x\cos^22x$.
\[\frac{\frac12\sin^22x\cos2x}{4\sin^22x\cos^22x}=\frac1{8\cos2x}.\]
\Javob{A}

\Savol{14}
$t=2^{\frac2x}>0$: \ $t^2-5t+4=0$ $\Rightarrow$ $t=1$ ($\frac2x=0$ — yöq) yoki $t=4$ $\Rightarrow$ $x=1$. \ $1$ ta.
\Javob{B}

\Savol{15}
$u=\log_4x$: \ $\log_x\sqrt{4x}=\frac12\left(\frac1u+1\right)$; öng tomon $-\frac1u\ge0$ $\Rightarrow$ $u<0$.
\[\frac12\left(\frac1u+1\right)=\frac1{u^2}\ \Rightarrow\ u^2+u-2=0\ \Rightarrow\ u=-2,\quad x=\frac1{16}.\]
\Javob{A}

\Savol{16}
$3x+6=\pm15$: \ $x=3$, $x=-7$. \ Köpaytma $-21$.
\Javob{A}

\Savol{17}
$x\ge0$, maxrajlar musbat. $x=0$ — ildiz ($\frac42=\frac42$).\\
$x>0$ da: $x^2-5x+4>0$ bölsa, çap tomon $>\dfrac{x^2-5x+4}{x+2\sqrt x+2}>$ öng tomon; $x^2-5x+4\le0$ bölsa, öng tomon $\le0<$ çap tomon. Yagona ildiz $0$.
\Javob{B}

\Savol{18}
$x\ge3$. $3\le x<6$ — hammasi yeçim. \ $x\ge6$: $4(x-3)>(x-6)^2$ $\Rightarrow$ $(x-4)(x-12)<0$ $\Rightarrow$ $x<12$.\\
$x\in[3;\,12)$: \ $3+4+\dots+11=63$.
\Javob{B}

\Savol{19}
$|x-3|\sqrt{(5-x)(x+2)}>0$ $\Rightarrow$ $-2<x<5$, $x\ne3$: \ $-1,\,0,\,1,\,2,\,4$ — $5$ ta.
\Javob{A}

\Savol{20}
$-2\le x\le5$; \ $x>0$, $x\ne1$; \ $\log_x2\ne1$ $\Rightarrow$ $x\ne2$; \ $\sin x\ne\frac\pi3$ har doim. \ Butun sonlar: $3,\,4,\,5$ — $3$ ta.
\Javob{D}

\Savol{21}
$t=\ln x$, $dt=\frac{dx}x$: \ $\displaystyle\int\ln t\,dt=t\ln t-t=\ln x\left(\ln\ln x-1\right)+C$.
\Javob{A}

\Savol{22}
$f'(x)=\frac32\sin^2\frac x2\cos\frac x2$: \ $f'\left(\frac{2\pi}3\right)=\frac32\cdot\frac34\cdot\frac12=\frac9{16}$.
\Javob{A}

\Savol{23}
Içki çizilgan aylana: $AB+CD=AD+BC$ $\Rightarrow$ $14=6+BC$, \ $BC=8$.
\Javob{B}

\Savol{24}
\fig{\begin{tikzpicture}[scale=0.32,font=\small]
\coordinate (B) at (-4,0); \coordinate (C) at (4,0); \coordinate (A) at (0,9.165); \coordinate (K) at (-2.4,3.666); \coordinate (M) at (0,0);
\draw (A)--(B)--(C)--cycle; \draw (0,2.619) circle (2.619); \fill (K) circle (3pt); \fill (M) circle (3pt); \draw[ax] (A)--(M);
\node[above] at (A) {$A$}; \node[below left] at (B) {$B$}; \node[below right] at (C) {$C$}; \node[below] at (M) {$M$};
\node[left] at (-1.3,6.6) {$6$}; \node[left] at (-3.3,1.8) {$4$}; \node[aux,below] at (-2,0) {$4$}; \node[aux,below] at (2,0) {$4$};
\end{tikzpicture}}
Urinma kesmalari: $B$ dan $4$ ($BC$ kiçik tomon), $A$ dan $6$. \ $AB=10$, $BM=4$, $BC=8$.\\
$AM=\sqrt{100-16}=2\sqrt{21}$: \ $S=\frac12\cdot8\cdot2\sqrt{21}=8\sqrt{21}$.
\Javob{D}

\Savol{25}
Kosinuslar teoremasi:
\[AB^2=450+25-2\cdot15\sqrt2\cdot5\cos135^\circ=475+150=625,\qquad AB=25.\]
\Javob{A}

\Savol{26}
\fig{\begin{tikzpicture}[scale=0.85,font=\small]
\coordinate (A) at (0,0); \coordinate (D) at (4,0); \coordinate (B) at (1,2.6); \coordinate (C) at (5,2.6); \coordinate (F) at (1,0); \coordinate (P) at (1.667,0.867); \coordinate (E) at (3,2.6);
\fill[okgreen!30] (P)--(F)--(D)--(C)--cycle;
\draw (A)--(B)--(C)--(D)--cycle; \draw (A)--(C); \draw (F)--(E);
\node[below left] at (A) {$A$}; \node[above left] at (B) {$B$}; \node[above right] at (C) {$C$}; \node[below] at (D) {$D$}; \node[below] at (F) {$F$}; \node[above] at (E) {$E$}; \node[above left] at (P) {$P$};
\end{tikzpicture}}
$\triangle APF\sim\triangle CPE$: $\dfrac{AF}{CE}=\dfrac{AP}{PC}=\dfrac12$ $\Rightarrow$ $CE=2AF=\dfrac12BC$. \ $S_{ABCD}=S$.
\[S_{PEC}=\frac{PC}{AC}\cdot\frac{CE}{CB}\cdot S_{ABC}=\frac23\cdot\frac12\cdot\frac S2=\frac S6,\qquad S_{ABEP}=\frac S2-\frac S6=\frac S3=16\ \Rightarrow\ S=48.\]
\[S_{APF}=\frac13\cdot\frac14\cdot\frac S2=2,\qquad S_{PFDC}=S_{ACD}-S_{APF}=24-2=22.\]
\Javob{A}

\Savol{27}
$t=\sqrt{9+x^2}\ge3$: \ $y=t^2-t-9$ — $t\ge3$ da ösuvçi, $y_{\min}=-3$.
\Javob{A}

\Savol{28}
Sakkizburçak içki burçagi $135^\circ$, beşburçakniki $108^\circ$; $AI$ — umumiy tomon:
\[\angle BAG=\angle BAI-\angle GAI=135^\circ-108^\circ=27^\circ.\]
\Javob{A}

\Savol{29}
$B(5;0)$, $D(0;4)$. $\overrightarrow{BE}=2\overrightarrow{BC}$ $\Rightarrow$ $E(5;8)$.\\
$ABED$ — trapetsiya: asoslari $AD=4$, $BE=8$, balandligi $AB=5$: \ $S=\frac{4+8}2\cdot5=30$.
\Javob{C}

\Savol{30}
\fig{\begin{tikzpicture}[scale=0.5,font=\small]
\coordinate (H) at (0,0); \coordinate (A) at (0,6); \coordinate (M) at (8,0); \coordinate (P) at (2.88,3.84);
\draw (H)--(M)--(A)--cycle; \draw[axs] (H)--(P);
\draw ($(H)!0.5cm!(M)$)--($(H)!0.5cm!(M)!0.5cm!90:(M)$)--($(H)!0.5cm!(A)$);
\node[below left] at (H) {$H$}; \node[above] at (A) {$A$}; \node[below right] at (M) {$M$}; \node[aux,above right] at (P) {$P$};
\node[left] at (0,3) {$6$}; \node at (6.6,0.4) {\scriptsize$37^\circ$};
\end{tikzpicture}}
$\angle HAM=53^\circ$: \ $HP=AH\sin53^\circ=AH\cos37^\circ=6\cdot0{,}8=4{,}8$.
\Javob{C}

\Savol{31}
$A\cup C=\{-4,\dots,5\}$, \ $B\cup D=\{-5,\dots,6\}$. \ $(A\cup C)'\cap(B\cup D)=\{-5;\,6\}$ — $2^2=4$ ta qism töplam.
\Javob{D}

\Savol{32}
Hammasi $9000$ ta. Oxirgi raqam $5$, qolgan uçtasi toq: $5^3=125$. \ $P=\dfrac{125}{9000}=\dfrac1{72}$.
\Javob{A}

\newpage
\Savol{33}
\fig{\begin{tikzpicture}[scale=0.14,font=\scriptsize]
\fill[okgreen!30] (10,10) .. controls (14,9) and (19,4) .. (20,0) .. controls (16,1) and (11,6) .. (10,10);
\draw (0,5) ellipse (1.6 and 5); \draw (0,10)--(10,10); \draw (0,0)--(20,0);
\draw (10,10)--(10,20); \draw (20,0)--(20,20); \draw (15,20) ellipse (5 and 1.6);
\draw (10,10) .. controls (14,9) and (19,4) .. (20,0); \draw[dashed] (10,10) .. controls (11,6) and (16,1) .. (20,0);
\draw[axs] (10,10)--(20,0);
\node[left] at (10,10) {$P$}; \node[below right] at (20,0) {$Q$};
\draw[{Stealth[length=3pt]}-{Stealth[length=3pt]}] (0,-2)--(20,-2); \node[below] at (10,-2) {$20$};
\draw[{Stealth[length=3pt]}-{Stealth[length=3pt]}] (0,12)--(10,12); \node[above] at (5,12) {$10$};
\end{tikzpicture}}
Har bir bölak — radiusi $r=5$ bölgan silindrning $45^\circ$ li tekislik bilan kesilgan qismi; yasovçilari $10$ dan $20$ gaça, örtaça uzunligi $15$ (bölakni xuddi şunday ikkinçi bölak bilan töldirsak, uzunligi $30$ bölgan silindr hosil böladi).

\textbf{33.} $S_{\text{yon}}=2\cdot2\pi r\cdot15=2\cdot2\cdot3\cdot5\cdot15=900$ — \textbf{B}.

\textbf{34.} $V=2\cdot\pi r^2\cdot15=2\cdot3\cdot25\cdot15=2250$ — \textbf{E}.

\textbf{35.} Kesim — ellips; uning katta öqi $PQ$ — tomoni $10$ bölgan kvadrat diagonali: $PQ=10\sqrt2$ — \textbf{C}.
\Javob{33 — B, \ 34 — E, \ 35 — C}

\newpage
\Savol{36}
Ikkinçi tenglama: $y^2-xy+2x-4=0$ $\Rightarrow$ $(y-2)(y+2)-x(y-2)=0$ $\Rightarrow$ $(y-2)(y+2-x)=0$.

$y=2$: \ $x^2-3x+9=9$ $\Rightarrow$ $x=0$ yoki $x=3$.

$x=y+2$: \ $(y+2)^2-3(y+2)+9=y^2+y+7=5y^2+y-13$ $\Rightarrow$ $y^2=5$, \ $y=\pm\sqrt5$, $x=2\pm\sqrt5$.

a) Yeçimlar: $(0;2)$, $(3;2)$, $\left(2+\sqrt5;\sqrt5\right)$, $\left(2-\sqrt5;-\sqrt5\right)$ — $4$ ta.

b) $0+3+\left(2+\sqrt5\right)+\left(2-\sqrt5\right)=7$.
\Javob{a) 4; \quad b) 7}

\newpage
\Savol{37}
$c=\cos3x$: \ $D=\left(4a^2-7\right)^2-8\left(2a^2-4\right)=\left(4a^2-9\right)^2$:
\[c=\frac{4a^2-7\pm\left(4a^2-9\right)}4\ \Rightarrow\ \cos3x=\frac12\ \ \text{yoki}\ \ \cos3x=2a^2-4.\]

a) $x\in\left[-\frac\pi2;\frac\pi2\right]$ $\Rightarrow$ $3x\in\left[-\frac{3\pi}2;\frac{3\pi}2\right]$: \ $\cos3x=\frac12$ dan $3x=\pm\frac\pi3$ — har doim $2$ ta ildiz. Eng kamida $2$ ta.

b) $3x=t\in\left[-\pi;\frac{3\pi}2\right]$. $\cos t=\frac12$: $t=\pm\frac\pi3$ — $2$ ta. Yana roppa-rosa bitta ildiz kerak.\\
$\cos t=m$ ning şu oraliqdagi ildizlari: $m=1$ da $1$ ta ($t=0$); $0<m<1$ da $2$ ta; $-1<m\le0$ da $3$ ta; $m=-1$ da $2$ ta.\\
Demak $2a^2-4=1$ $\Rightarrow$ $a^2=\frac52$, \ $a=\dfrac{\sqrt{10}}2$.
\Javob{a) 2; \quad b) $\dfrac{\sqrt{10}}2$}

\newpage
\Savol{38}
Simmetriya öqi $x=1$: \ $k=1+3=4$. \ $f(x)=a(x+2)(x-4)$, \ $f(1)=-9a=4$ $\Rightarrow$ $a=-\frac49$.

a) $a+b+c=f(1)=4$: \ $\dfrac4k=\dfrac44=1$.

b) $g(x)=3f\left(\frac x5\right)-2$: uç $\frac x5=1$ da, ya'ni $x=5$, \ $y=3\cdot4-2=10$. \ $5+10=15$.
\Javob{a) 1; \quad b) 15}

\newpage
\Savol{39}
\fig{\begin{tikzpicture}[x=0.5cm,y=0.6cm,font=\small]
\draw[-{Stealth[length=4pt]}] (-10,0)--(4.5,0) node[right] {$x$}; \draw[-{Stealth[length=4pt]}] (0,-3)--(0,2.2) node[left] {$y$};
\draw[thick] plot[smooth,tension=0.7] coordinates {(-9,-1.1) (-8.3,-2.3) (-7.6,-2.6) (-6.5,-1.8) (-5,0) (-4.2,1.1) (-3.6,0.6) (-3,0) (-2.3,-1.4) (-1,-2) (0.5,-1.8) (2,-1.2) (3,-0.5)};
\draw[fill=white] (-9,-1.1) circle (2.5pt); \draw[fill=white] (3,-0.5) circle (2.5pt);
\node[aux] at (-7,0.5) {$-$}; \node[aux] at (-4,-0.5) {$+$}; \node[aux] at (0.8,0.5) {$-$};
\fill (-5,0) circle (2pt); \fill (-3,0) circle (2pt);
\node[above left] at (-5,0) {$-5$}; \node[above right] at (-3,0) {$-3$};
\end{tikzpicture}}
$f'$ ning işorasi: $(-9;-5)$ da $-$, \ $(-5;-3)$ da $+$, \ $(-3;3)$ da $-$.

a) $(-5;3)$: $f$ avval ösadi, $x=-3$ dan keyin kamayadi — eng katta qiymat $x=-3$ da.

b) $(-9;-3)$: $f$ avval kamayadi, $x=-5$ dan keyin ösadi — eng kiçik qiymat $x=-5$ da.
\Javob{a) $-3$; \quad b) $-5$}

\newpage
\Savol{40}
$\sqrt{x-1}+2=\dfrac{x+3}2$ $\Rightarrow$ $\sqrt{x-1}=\dfrac{x-1}2$. \ $t=\sqrt{x-1}$: \ $t^2=2t$ $\Rightarrow$ $t=0$, $t=2$ $\Rightarrow$ $x=1$, $x=5$.\\
$A(1;\,2)$, \ $B(5;\,4)$.

a) $AB=\sqrt{4^2+2^2}=2\sqrt5$.

b) $[1;5]$ da $y_1\ge y_2$; \ $u=x-1$:
\[S=\int_1^5\left(\sqrt{x-1}-\frac{x-1}2\right)dx=\int_0^4\left(\sqrt u-\frac u2\right)du=\frac23\cdot8-4=\frac43.\]
\Javob{a) $2\sqrt5$; \quad b) $\dfrac43$}

\newpage
\Savol{41}
\fig{\begin{tikzpicture}[scale=0.62,font=\small]
\coordinate (A) at (0,0); \coordinate (B) at (0,5.657); \coordinate (C) at (4.243,0); \coordinate (D) at (0,4.243); \coordinate (F) at (2.121,2.121); \coordinate (E) at (3.536,0.707);
\fill[okgreen!30] (A)--(D)--(F)--cycle; \fill[okgreen!30] (B)--(E)--(C)--cycle;
\draw (A)--(B)--(C)--cycle; \draw (D)--(C); \draw (A)--(F); \draw (B)--(E);
\draw (0,0.35)--(0.35,0.35)--(0.35,0);
\draw ($(F)!0.3cm!(A)$)--($(F)!0.3cm!(A)!0.3cm!90:(A)$)--($(F)!0.3cm!(C)$);
\draw ($(D)!0.35cm!(B)$) arc [start angle=90,end angle=-45,radius=0.35cm]; \node at (0.5,4.85) {\scriptsize$\alpha$};
\node[below left] at (A) {$A$}; \node[left] at (B) {$B$}; \node[right] at (C) {$C$}; \node[left] at (D) {$D$}; \node[above right] at (F) {$F$}; \node[below left] at (E) {$E$};
\node[aux,left] at (0,4.95) {$\sqrt2$}; \node[aux] at (1.3,0.85) {$3$};
\end{tikzpicture}}
$\angle ADC=180^\circ-135^\circ=45^\circ$ $\Rightarrow$ $\triangle ADC$ — teng yonli töğri burçakli, $AF$ — balandlik va mediana:
\[DF=FC=AF=3,\quad AD=AC=3\sqrt2,\quad AB=AD+BD=4\sqrt2.\]

a) $BC=\sqrt{32+18}=5\sqrt2$.

b) $S_{ADF}=\frac12\cdot3\cdot3=\frac92$. \ $S_{BDC}=\frac12\cdot BD\cdot AC=\frac12\cdot\sqrt2\cdot3\sqrt2=3$, \ $EC=1=\frac16DC$:
\[S_{BEC}=\frac16\cdot3=\frac12,\qquad S=\frac92+\frac12=5.\]
\Javob{a) $5\sqrt2$; \quad b) 5}

\newpage
\Savol{42}
\fig{\begin{tikzpicture}[scale=0.55,font=\small]
\coordinate (A) at (0,0); \coordinate (D) at (8,0); \coordinate (B) at (2.5,4.33); \coordinate (C) at (10.5,4.33); \coordinate (O) at (3,1.732); \coordinate (T) at (3,0);
\draw (A)--(B)--(C)--(D)--cycle; \draw (B)--(D); \draw (O) circle (1.732); \draw (O)--(T); \draw[ax] (A)--(O); \fill (O) circle (2pt);
\node[below left] at (A) {$A$}; \node[above] at (B) {$B$}; \node[above right] at (C) {$C$}; \node[below] at (D) {$D$}; \node[left] at (2.9,2) {$O$}; \node[right] at (3,0.8) {$\sqrt3$};
\node[aux,below] at (T) {$T$}; \node at (1.05,0.3) {\scriptsize$30^\circ$};
\end{tikzpicture}}
$\angle A=60^\circ$, \ $AB+AD=13$. Urinma kesma: $AT=r\ctg30^\circ=\sqrt3\cdot\sqrt3=3$.\\
$AT=p-BD$, \ $p=\dfrac{13+BD}2$ $\Rightarrow$ $3=\dfrac{13-BD}2$, \ $BD=7$.\\
Kosinuslar teoremasi: $49=AB^2+AD^2-AB\cdot AD=169-3\,AB\cdot AD$ $\Rightarrow$ $AB\cdot AD=40$.

a) $S=AB\cdot AD\sin60^\circ=40\cdot\dfrac{\sqrt3}2=20\sqrt3$.

b) $AB+AD=13$, $AB\cdot AD=40$ $\Rightarrow$ $5$ va $8$. Katta tomon $8$.
\Javob{a) $20\sqrt3$; \quad b) 8}

\newpage
\Savol{43}
\fig{\begin{tikzpicture}[scale=1.6,font=\small]
\foreach \n/\a in {A/105,B/75,C/45,D/15,E/-15,F/-45,G/-75,H/-105,I/-135,J/-165,K/165,L/135} {\coordinate (\n) at (\a:1); \node at (\a:1.15) {$\n$};}
\fill[okgreen!30] (L)--(E)--(F)--cycle;
\draw (A)--(B)--(C)--(D)--(E)--(F)--(G)--(H)--(I)--(J)--(K)--(L)--cycle; \draw (L)--(C); \draw (L)--(E); \draw (L)--(F);
\coordinate (O) at (0,0); \fill (O) circle (0.8pt); \node[below] at (O) {$O$}; \draw[ax] (O)--(L); \draw[ax] (O)--(C);
\end{tikzpicture}}
Markaziy burçak $30^\circ$. \ $a=2R\sin15^\circ=R\cdot\dfrac{\sqrt6-\sqrt2}2$ $\Rightarrow$ $R=\dfrac{2\left(\sqrt3-1\right)}{\sqrt2\left(\sqrt3-1\right)}=\sqrt2$.

a) $L$ dan $C$ gaça $3$ ta tomon: $\angle LOC=90^\circ$, \ $LC=R\sqrt2=2$.

b) $LF$ — diametr ($6$ ta tomon) $\Rightarrow$ $\angle LEF=90^\circ$. \ $LF=2\sqrt2$, $EF=\sqrt3-1$:
\[LE=\sqrt{8-\left(4-2\sqrt3\right)}=\sqrt{4+2\sqrt3}=\sqrt3+1,\qquad S=\frac12\left(\sqrt3+1\right)\left(\sqrt3-1\right)=1.\]
\Javob{a) 2; \quad b) 1}

\newpage
\Savol{44}
\fig{\begin{tikzpicture}[scale=1,font=\small]
\coordinate (O) at (0,2.6); \coordinate (C) at (-1.5,0); \coordinate (Q) at (1.5,0); \coordinate (D) at (0,0);
\draw (C)--(O)--(Q); \draw (C) arc (180:360:1.5 and 0.3); \draw[dashed] (Q) arc (0:180:1.5 and 0.3);
\draw[axs] (O)--(D)--(Q); \draw[green!50!black] (0,0.866) circle (0.866); \fill (0,0.866) circle (1.2pt);
\node[above] at (O) {$O$}; \node[left] at (C) {$C$}; \node[right] at (Q) {$A,\,B$}; \node[below] at (D) {$D$};
\node[right] at (0.8,1.4) {$R$}; \node[aux,below] at (0.75,0) {$r$}; \node at (0.25,2.0) {\scriptsize$\varphi$};
\end{tikzpicture}}
Konus yasovçisi $l=R$; asos aylanasi uzunligi yarim aylana yoyiga teng: $2\pi r=\pi R$ $\Rightarrow$ $r=\dfrac R2$.

a) $\sin\varphi=\dfrac rl=\dfrac12$.

b) $\varphi=30^\circ$ $\Rightarrow$ öq kesimi — tomoni $2r=R=2\sqrt3$ bölgan teng tomonli uçburçak.\\
Içki çizilgan sfera radiusi — şu uçburçak içki aylanasi radiusi: $\rho=\dfrac{2\sqrt3}{2\sqrt3}=1$.
\[S=4\pi\rho^2=4\pi=12.\]
\Javob{a) $\dfrac12$; \quad b) 12}

\newpage
\Savol{45}
$10$ km yölga $t=\dfrac{10}v$ soat ketadi. Harajat:
\[C(v)=\frac{10}v\left(250+v^3\right)=\frac{2500}v+10v^2,\qquad C'(v)=-\frac{2500}{v^2}+20v.\]
$C'(v)=0$ $\Rightarrow$ $v^3=125$, \ $v=5$ ($v<5$ da $C'<0$, $v>5$ da $C'>0$ — minimum).

a) $v=5$ km/h.

b) $C(5)=\dfrac{2500}5+10\cdot25=500+250=750$ \$.
\Javob{a) 5; \quad b) 750}

\ShowJavoblar
\end{document}
